\section{Rank-\texorpdfstring{$l$}{l} transport of the original relation-Gram ratio}

This calculation keeps the two fixed polynomial inclusions and transports
their source form through the same positive density. The difference of
their orthogonal projectors has precisely $l$ positive and $l$ negative
eigenvalues. Its signed derivative controls the complete Gram-determinant
ratio by the degree of the original multiplication polynomial.

\subsection{Fixed coefficient columns, forms, and exact intersection}
Let $0<\epsilon<64$, put
$\mathcal K_\epsilon=[-64,-\epsilon]\cup[\epsilon,64]$, and retain the
evaluation coordinate $S=c+iQz$, where $c\in\mathbb R$, $Q>0$.
Use a finite positive Borel measure $w$ on $\mathcal K_\epsilon$ whose
support is infinite. This is a specified domain of the Gram construction;
the concrete original densities used below are strictly positive on both
intervals, and their finiteness and support are verified below.
Let $\psi$ be a bounded real measurable function, $w_t=e^{t\psi}w$ for
$0\leq t\leq1$, and retain a monic polynomial
\[
 f(S)=\sum_{j=0}^l f_jS^j,\qquad f_l=1,\quad
 l\geq1,\quad r\geq l,\quad n=r+l.
 \tag{CTR1}
\]
The abstract polynomial spaces and their coefficient columns are
\[
\begin{gathered}
 \mathcal H=\mathbb C[S]_{\leq r+l-1},\qquad
 \mathcal P=\mathbb C[S]_{\leq r-1},\qquad
 \mathcal F=f\mathcal P,\\
 C_0=[e_0,\ldots,e_{r-1}],\quad
 C_f=[\operatorname{coeff}(fS^j)]_{j=0}^{r-1}.
\end{gathered}
 \tag{CTR2}
\]
Here $e_j$ is the coefficient column of $S^j$ in the original ordered
ambient basis $1,S,\ldots,S^{n-1}$. Thus $C_0,C_f$ are fixed $n\times r$
matrices, and $(C_f)_{a,j}=f_{a-j}$, where a coefficient outside $0,\ldots,l$
is zero. Multiplication by a nonzero polynomial is injective, so both
matrices have rank $r$.

The original source form, its two restriction Grams, and the full ratio are
\[
 \begin{split}
 (M_t)_{a,b}&=\int_{\mathcal K_\epsilon}
        \overline{(c+iQz)^a}(c+iQz)^b e^{t\psi(z)}\,dw(z),\quad 0\leq a,b<n,\\
 G_0(t)&=C_0^*M_tC_0,\qquad G_f(t)=C_f^*M_tC_f,\\
 \mathcal R_{w_t}(f)&=\frac{\det G_f(t)}{\det G_0(t)}.
 \end{split}
 \tag{CTR3}
\]
Every entry is finite because the coordinate and $\psi$ are bounded on
the compact set and $w$ is finite. Every $M_t$ is positive definite:
a nonzero polynomial has finitely many zeros on the injective line
$z\mapsto c+iQz$, whereas the support of $w$ is infinite; at a support
point outside those zeros it is bounded away from zero on a neighbourhood
of positive measure. Multiplication by the positive bounded density
$e^{t\psi}$ preserves that conclusion. Both restriction Grams are
therefore positive definite.

There is an exact sequence, including all coefficient maps,
\[
 0\longrightarrow\mathbb C[S]_{\leq r-l-1}
 \xrightarrow{\ a\mapsto(-fa,a)\ }
 \mathcal P\oplus\mathcal P
 \xrightarrow{\ (u,v)\mapsto u+fv\ }
 \mathcal H\longrightarrow0.
 \tag{CTR4}
\]
The convention $\mathbb C[S]_{\leq-1}=0$ applies at $r=l$.
Indeed, division of $h\in\mathcal H$ by the original monic $f$ gives
$h=fv+u$ with $\deg v\leq r-1$ and $\deg u<l\leq r$. The kernel is
$u=-fv$, and its degree constraint forces $\deg v\leq r-l-1$.
This proves both surjectivity and the kernel identity. In particular,
\[
 \mathcal I:=\mathcal P\cap\mathcal F
       =f\mathbb C[S]_{\leq r-l-1},\quad
 \dim\mathcal I=r-l,\qquad \mathcal P+\mathcal F=\mathcal H.
 \tag{CTR5}
\]
No source metric was needed for this exact polynomial intersection.

For any one of the fixed injective column matrices $C$ define
\[
 P_C(t)=C(C^*M_tC)^{-1}C^*M_t.
 \tag{CTR6}
\]
Direct multiplication gives $P_C^2=P_C$, $P_CC=C$ and
$P_C^*M_t=M_tP_C$. Also $C^*M_t(x-P_Cx)=0$, proving that this is the
actual $M_t$-orthogonal projector onto the column space of $C$.
Cyclicity gives $\operatorname{Tr}P_C=\operatorname{Tr}I_{\operatorname{rank}C}$.
Let $P_0=P_{C_0}$, $P_f=P_{C_f}$ and
\[
 D_t=P_f(t)-P_0(t),\qquad
 D_t^*M_t=M_tD_t,\qquad \operatorname{Tr}D_t=r-r=0.
 \tag{CTR7}
\]

\subsection{The complete residual frame and its signed spectrum}
The columns $C_I=[\operatorname{coeff}(fS^j)]_{j=0}^{r-l-1}$ give the
intersection, and $P_I$ denotes CTR6 for $C_I$; if $r=l$, $P_I=0$.
The inclusions $\mathcal I\subset\mathcal P,\mathcal F$ imply
$P_0P_I=P_fP_I=P_I$ and, by the adjoint equations,
$P_IP_0=P_IP_f=P_I$. Define the following literal columns:
\[
 \begin{split}
 U_0&=(I-P_I)[e_0,\ldots,e_{l-1}],\\
 U_f&=(I-P_I)[\operatorname{coeff}(fS^{r-l+j})]_{j=0}^{l-1},\\
 H_0&=U_0^*M_tU_0,\qquad H_f=U_f^*M_tU_f,\qquad C=U_0^*M_tU_f.
 \end{split}
 \tag{CTR8}
\]
Each $U$ has rank $l$. For $U_0$, a polynomial of degree less than $l$
belonging to $\mathcal I$ must be zero. For $U_f$, divisibility by $f$
reduces a dependence in $\mathcal I$ to a polynomial supported in degrees
$r-l,\ldots,r-1$ and at the same time of degree at most $r-l-1$.
Their ranges are $\mathcal P\cap\mathcal I^{\perp_{M_t}}$ and
$\mathcal F\cap\mathcal I^{\perp_{M_t}}$, respectively. Thus
\[
 P_0=P_I+U_0H_0^{-1}U_0^*M_t,\qquad
 P_f=P_I+U_fH_f^{-1}U_f^*M_t.
 \tag{CTR9}
\]
These are orthogonal decompositions inside each individual space; the
cross Gram $C$ between the two residual spaces remains present.
By CTR5, the frame $B=[C_I,U_0,U_f]$ is invertible and satisfies
\[
 B^*M_tB=
 \begin{pmatrix}C_I^*M_tC_I&0&0\\0&H_0&C\\0&C^*&H_f\end{pmatrix},\qquad
 B^{-1}D_tB=
 \begin{pmatrix}0&0&0\\0&-I_l&-H_0^{-1}C\\0&H_f^{-1}C^*&I_l\end{pmatrix}.
 \tag{CTR10}
\]
Both formulas follow by applying CTR9 separately to its three displayed
column groups; the second keeps both off-diagonal blocks with their signs.

Here is an exact spectral frame for this same endomorphism. Take the
positive roots of $H_0,H_f$ and put
$T=H_0^{-1/2}CH_f^{-1/2}$. Positive definiteness of the lower Gram in
CTR10 implies $I-T^*T>0$, by its Schur complement. Choose orthonormal
singular vectors $a_j,b_j$ with
$a_i^*Tb_j=\sigma_j\delta_{ij}$ and $0\leq\sigma_j<1$.
They can be constructed by diagonalizing the Hermitian matrix $T^*T$:
for a positive eigenvalue $\sigma_j^2$, set $a_j=Tb_j/\sigma_j$;
complete the resulting left vectors by an orthonormal basis of
$\ker T^*$ at the zero singular values. This gives both orthonormal
bases and the stated equation. The finite Hermitian spectral theorem
itself follows by choosing one eigenvector, noting that its eigenvalue
is real and its orthogonal complement is invariant, and then inducting
on the dimension. This also constructs the positive square roots above.

Set, in the original coefficient space,
\[
 u_j=U_0H_0^{-1/2}a_j,\qquad v_j=U_fH_f^{-1/2}b_j,\qquad
 s_j=(1-\sigma_j^2)^{1/2}>0,\qquad
 w_j=\frac{v_j-\sigma_ju_j}{s_j}.
 \tag{CTR11}
\]
Their actual inner products are
$u_i^*M_tu_j=v_i^*M_tv_j=\delta_{ij}$ and
$u_i^*M_tv_j=\sigma_j\delta_{ij}$; hence $(u_j,w_j)_{j=1}^l$ is an
$M_t$-orthonormal frame of $\mathcal I^{\perp_{M_t}}$.
On its $j$th two-dimensional plane, CTR9 acts by
\[
 P_0=\begin{pmatrix}1&0\\0&0\end{pmatrix},\qquad
 P_f=\begin{pmatrix}\sigma_j^2&\sigma_js_j\\\sigma_js_j&s_j^2\end{pmatrix},\qquad
 D_t=\begin{pmatrix}-s_j^2&\sigma_js_j\\\sigma_js_j&s_j^2\end{pmatrix}.
 \tag{CTR12}
\]
Write $\alpha_j=\sqrt{(1-s_j)/2}$ and
$\beta_j=\sqrt{(1+s_j)/2}$. Direct multiplication gives the orthonormal
eigenvectors
\[
 e_j^+=\alpha_ju_j+\beta_jw_j,\qquad
 e_j^-=\beta_ju_j-\alpha_jw_j,\qquad
 D_te_j^+=s_je_j^+,\quad D_te_j^-=-s_je_j^-.
 \tag{CTR13}
\]
This includes $\sigma_j=0$, when $s_j=1$; no division by $\sigma_j$
is used in CTR11 or CTR13. On $\mathcal I$ both projectors are the
identity and their difference is zero. Consequently the complete inertia,
rank, and trace of the original-metric absolute value are
\[
 \operatorname{inertia}(D_t)=(l,l,r-l),\qquad
 \operatorname{rank}D_t=2l,\qquad
 \operatorname{spec}D_t\subset[-1,1],\qquad
 \operatorname{Tr}|D_t|=2\sum_{j=1}^l s_j\leq2l.
 \tag{CTR14}
\]
Here $|D_t|$ is defined by the displayed $M_t$-orthonormal eigenframe;
equivalently it is the positive root of $D_t^2$ in that same metric.
These are the signed spectral maps of the original two subspaces,
with their actual cross Gram, rather than an arbitrary rank-$2l$ matrix.

\subsection{The derivative of the full ratio and its integrated bound}
Let $a=\operatorname*{ess\,inf}_w\psi$,
$b=\operatorname*{ess\,sup}_w\psi$, and $\operatorname{osc}_w\psi=b-a$.
For $t\in[0,1]$ differentiation under each finite Gram integral gives
\[
 (M_t')_{i,j}=\int_{\mathcal K_\epsilon}
       \overline{(c+iQz)^i}(c+iQz)^j\psi(z)e^{t\psi(z)}\,dw(z),\qquad
 A_t=M_t^{-1}M_t'.
 \tag{CTR15}
\]
To justify differentiation, extend $t$ to $[-1,2]$. The derivative
integrand is bounded in absolute value by a fixed constant times
$\|\psi\|_\infty e^{2\|\psi\|_\infty}$, integrable because $w$ is
finite and the polynomial coordinates are bounded. The difference
quotients have the same domination by the mean value theorem.
Dominated convergence proves the derivative and its continuity.

The matrix $A_t$ is the compression of multiplication by $\psi$ to
the original polynomial image in $L^2(w_t)$. More precisely, if
$\mathcal E:\mathbb C^n\to L^2(w_t)$ is coefficient evaluation, then
$\mathcal E^*\mathcal E=M_t$,
$\mathcal E^*M_\psi\mathcal E=M_t'$, and its coefficient compression is
$M_t^{-1}\mathcal E^*M_\psi\mathcal E=A_t$.
It satisfies $A_t^*M_t=M_tA_t=M_t'$ and, for every coefficient $x$,
\[
 a\,x^*M_tx\leq x^*M_tA_tx
   =\int\psi(z)|\mathcal Ex(z)|^2\,dw_t(z)
 \leq b\,x^*M_tx.
 \tag{CTR16}
\]
The notation $aI\leq A_t\leq bI$ means exactly this original-metric
quadratic-form inequality.

Jacobi's determinant formula, with the fixed columns in CTR2, now gives
\[
 \begin{split}
 \frac{d}{dt}\log\mathcal R_{w_t}(f)
 &=\operatorname{Tr}\bigl(G_f^{-1}C_f^*M_t'C_f\bigr)
   -\operatorname{Tr}\bigl(G_0^{-1}C_0^*M_t'C_0\bigr)\\
 &=\operatorname{Tr}\bigl((P_f-P_0)M_t^{-1}M_t'\bigr)
   =\operatorname{Tr}(D_tA_t).
 \end{split}
 \tag{CTR17}
\]
For completeness, Jacobi's formula follows by multilinearity of the
determinant, differentiating one column at a time, and then using
$\operatorname{adj}(G)=\det(G)G^{-1}$; all $G$ here are invertible by
CTR3. Cyclicity and $M_tM_t^{-1}=I$ prove the second equality in CTR17.

In the complete actual eigenframe CTR13,
\[
 \operatorname{Tr}(D_tA_t)
 =\sum_{j=1}^l s_j
   \bigl((e_j^+)^*M_tA_te_j^+-(e_j^-)^*M_tA_te_j^-\bigr).
 \tag{CTR18}
\]
The intersection contributes zero. Each scalar in parentheses is the
difference of two numbers in $[a,b]$ by CTR16; therefore
\[
 |\operatorname{Tr}(D_tA_t)|
 \leq\frac{\operatorname{Tr}|D_t|}{2}\,(b-a)
 \leq l\operatorname{osc}_w\psi.
 \tag{CTR19}
\]
This also follows by centering $A_t$ at $(a+b)I/2$ and using
$\operatorname{Tr}D_t=0$; the signed scalar cancellation is exact.
All $G_t$ and their inverses vary continuously on the compact $t$
interval, so CTR17 is a continuous scalar derivative. Integration yields
the completed logarithmic transport estimate
\[
 \boxed{
 \left|\log\mathcal R_{e^\psi w}(f)-\log\mathcal R_w(f)\right|
 \leq l\operatorname{osc}_w\psi
 \leq2l\|\psi\|_\infty.}
 \tag{CTR20}
\]
If $l=0$ and $f=1$, the two inclusions coincide and their ratio is
identically one, giving the zero right side directly.

\subsection{Application to the full original arithmetic polynomial}
Retain the original packet, including every multiplicity:
\[
 \begin{split}
 k&=4l+1,\quad l\geq1,\quad m\geq1,\quad e=1+k(m-1),\\
 q&=e(k+1)^2,\quad c=k/2,\quad r\in\{q,q+1\},\\
 d&=\delta^2,\quad g=\gamma^2,\quad 0<\delta<\tfrac12,\quad\gamma>2,\\
 \chi(c+x)&=\prod_{a,b=0}^k
       [x-(2a-k)\delta-i(2b-k)\gamma]^e,\qquad
 f(S)=\prod_{j=0}^{l-1}(S-c-2j-\tfrac12).
 \end{split}
 \tag{CTR21}
\]
Use $Q=q$, $\epsilon=2^{-10}$ and the same compact bulk
$\mathcal K_\epsilon$. The concrete positive base measure, with its full
source mass and Jacobian, is
\[
 dw_{\mathrm{pow}}(z)
   =q^{2q+1}|z|^{2q}\sigma(qz)\,dz,\qquad
 \sigma(y)=\frac{|\Gamma(1/4+iy/2)|^2}{2\pi}.
 \tag{CTR22}
\]
For each fixed original $q$, this is a finite measure with strictly
positive continuous density on both compact intervals: the Gamma
argument has positive real part, so the original Gamma factor is finite,
continuous, and nonzero there. Thus every Gram required in CTR3 is
positive definite for this actual source; no moment or support
nondegeneracy is assumed without verification.

Put
\[
 \begin{split}
 a_k&=\frac{k(k+2)(d-g)}{3q},\qquad
 \tau_\epsilon=\frac{k^2(d+g)}{q^2\epsilon^2},\qquad
 \delta_\chi=\frac{q\tau_\epsilon^2}{2(1-\tau_\epsilon)},\\
 \psi_\chi(z)&=\log\frac{|\chi(c+iqz)|^2}{(q|z|)^{2q}}.
 \end{split}
 \tag{CTR23}
\]
The finite condition $k\geq2048\sqrt{d+g}$ ensures
$k\sqrt{d+g}/q\leq\epsilon/2$ and hence $\tau_\epsilon\leq1/4$,
because $q\geq k^2$ and $2/\epsilon=2048$.
It makes every factor of $\chi(c+iqz)$ nonzero on the bulk, so the
logarithm in CTR23 is a bounded continuous real function there.

Here is the exact quadratic expansion needed for its oscillation.
Set $\zeta_{ab}=(2a-k)\delta+i(2b-k)\gamma$.
For $|y|>k\sqrt{d+g}$,
\[
 \log\frac{|\chi(c+iy)|^2}{|y|^{2q}}
 =2e\operatorname{Re}\sum_{a,b}\log(1+i\zeta_{ab}/y)
 =\frac{p_2}{y^2}+E_\chi(y),\qquad
 p_2=\frac{qk(k+2)}3(d-g).
 \tag{CTR24}
\]
The Taylor series of $\log(1+u)$ is absolutely convergent at every
displayed factor. Odd powers cancel under the exact permutation
$(a,b)\mapsto(k-a,k-b)$. The second power is
$e\operatorname{Re}\sum\zeta_{ab}^2/y^2$.
The elementary finite sums
$\sum_{a=0}^k(2a-k)=0$ and
$\sum_{a=0}^k(2a-k)^2=(k+1)k(k+2)/3$
follow by expanding and using
$\sum a=k(k+1)/2$ and $\sum a^2=k(k+1)(2k+1)/6$.
They give the coefficient $p_2$ with its sign and multiplicity in CTR24.
For the remaining even terms put $\tau=k^2(d+g)/y^2<1$.
There are exactly $q$ roots counted with multiplicity, so
\[
 |E_\chi(y)|
 \leq q\sum_{j=2}^\infty\frac{\tau^j}{j}
 \leq\frac{q\tau^2}{2(1-\tau)}.
 \tag{CTR25}
\]
The first inequality bounds each real root power by its modulus and
keeps the cancellation of odd powers already proved; the second uses
$1/j\leq1/2$ for $j\geq2$ and sums the geometric series.
Substitution $y=qz$ therefore proves on the whole bulk
\[
 \psi_\chi(z)=\frac{a_k}{z^2}+E_\chi(qz),\qquad
 |E_\chi(qz)|\leq\delta_\chi,\qquad
 \operatorname{osc}\psi_\chi
 \leq|a_k|(\epsilon^{-2}-64^{-2})+2\delta_\chi
 \leq2(|a_k|\epsilon^{-2}+\delta_\chi).
 \tag{CTR26}
\]
For the first oscillation bound, $z^{-2}$ lies in
$[64^{-2},\epsilon^{-2}]$ and the remainder's range lies in
$[-\delta_\chi,\delta_\chi]$. This retains the quadratic signed term
before bounding its oscillation.

The actual and quadratic-deformed measures are
\[
 \begin{split}
 dw_\chi&=q|\chi(c+iqz)|^2\sigma(qz)\,dz=e^{\psi_\chi}dw_{\mathrm{pow}},\\
 dw_{\mathrm{\quad}}&=e^{a_k/z^2}dw_{\mathrm{pow}},\qquad
 dw_\chi=e^{E_\chi(qz)}dw_{\mathrm{\quad}}.
 \end{split}
 \tag{CTR27}
\]
Each measure has the verified positivity and finiteness of CTR22.
Apply CTR20 with the same $\mathcal P$, same original $f$, and the
same coordinate map $P(S)\mapsto P(c+iqz)$, to obtain
\[
 \boxed{
 \begin{aligned}
 |\log\mathcal R_{w_\chi}(f)-\log\mathcal R_{w_{\mathrm{pow}}}(f)|
 &\leq l\{|a_k|(\epsilon^{-2}-64^{-2})+2\delta_\chi\}\\
 &\leq2l(|a_k|\epsilon^{-2}+\delta_\chi),\\
 |\log\mathcal R_{w_\chi}(f)-\log\mathcal R_{w_{\mathrm{\quad}}}(f)|
 &\leq2l\delta_\chi.
 \end{aligned}}
 \tag{CTR28}
\]
Both statements hold for $r=q$ and $r=q+1$. Their coefficient is $l$,
although each individual Gram determinant has $r$ columns.

These are two forms on the same original relation coefficient space.
Their exact realization inside the one bulk Hilbert space
$L^2(\mathcal K_\epsilon,q\sigma(qz)dz)$ is
\[
 \begin{split}
 T_{\mathrm{pow}}P(z)&=(iqz)^qP(c+iqz),\qquad
 T_\chi P(z)=\chi(c+iqz)P(c+iqz),\\
 (Uh)(z)&=\frac{\chi(c+iqz)}{(iqz)^q}h(z),\qquad
 UT_{\mathrm{pow}}=T_\chi,\qquad
 UM_{f(c+iqz)}=M_{f(c+iqz)}U.
 \end{split}
 \tag{CTR29}
\]
The multiplier and its inverse are bounded on the compact bulk because
their denominators and numerators are nonzero there. Its squared modulus
is exactly $e^{\psi_\chi}$. Taking the Hilbert Gram of these maps yields
CTR27, including the numerator maps obtained by composing with the same
$f$. Thus CTR29 proves the source-form comparison without introducing
any replacement for the original arithmetic quotient
$E_\chi=\mathbb C[S]/(\chi)$.

If monic $z$ coordinates are used later, the exact polynomial is
$\widetilde f(z)=(iq)^{-l}f(c+iqz)$, and the coefficient change
$P(S)\mapsto P(c+iqz)$ is triangular with diagonal $(iq)^j$,
$0\leq j<r$. Its common squared determinant cancels between the two
restriction Grams, whereas the leading factor of $f$ gives
\[
 \mathcal R_w^{S}(f)=q^{2lr}\mathcal R_w^{z}(\widetilde f),\qquad
 \log\mathcal R_w^{S}(f)=2lr\log q+\log\mathcal R_w^{z}(\widetilde f).
 \tag{CTR30}
\]
Indeed each of the $r$ numerator polynomial columns acquires the factor
$(iq)^l$, and its Gram determinant acquires $|(iq)^l|^{2r}=q^{2lr}$.
This factor is retained separately in each source form before their
logarithmic difference is taken.

Finally the bounds in CTR28 have explicit original-degree rates.
On the displayed threshold $\tau_\epsilon\leq1/4$,
\[
 |a_k|\leq\frac{|d-g|}{3e},\qquad
 \delta_\chi\leq
 \frac{2k^4(d+g)^2}{3e^3(k+1)^6\epsilon^4}
 \leq\frac{2(d+g)^2}{3e^3k^2\epsilon^4}.
 \tag{CTR31}
\]
The first inequality uses $k(k+2)<(k+1)^2$. The second substitutes the
original $q=e(k+1)^2$ into CTR23 and uses $1/(1-\tau_\epsilon)\leq4/3$.
Since $2l=(k-1)/2$, the coarser first error in CTR28 is bounded by
\[
 \frac{(k-1)|d-g|}{6e\epsilon^2}
 +\frac{(k-1)(d+g)^2}{3e^3k^2\epsilon^4}.
 \tag{CTR32}
\]
For fixed $\delta,\gamma,m$ and fixed $\epsilon=2^{-10}$ this is
$O(k/e)$, hence $O(k)=o(q)$ because $q\geq(k+1)^2$.
For fixed $m\geq2$, $e\geq k(m-1)$, so the same explicit bound is
even $O(1)$. The quadratic-retaining second error is bounded by
\[
 2l\delta_\chi\leq
 \frac{(k-1)(d+g)^2}{3e^3k^2\epsilon^4},
 \tag{CTR33}
\]
which is $O(k^{-1})$ at $m=1$ and $O(k^{-4})$ for every fixed $m\geq2$.
These are bounds on the complete restricted relation-Gram ratio with
its original multiplication by $f$. The lower and far-tail source
contributions remain in their separate full-source comparison; no
bulk estimate is substituted for those omitted integrals.

\subsection{Complete return to the original sigma shifted-moment determinants}
We now transport the entire line integral, using the proved original
inner and far estimates and retaining the exact four tail corrections.
For $s=0,l$ and $0\leq i,j<r$, define the full real moment matrices
\[
 (\mathsf M_{q+s,r})_{i,j}
   =\int_{\mathbb R}y^{2(q+s)+i+j}\sigma(y)\,dy,\qquad
 \mathcal H_{q,l,r}
   =\frac{\det\mathsf M_{q+l,r}}{\det\mathsf M_{q,r}}.
 \tag{CTR34}
\]
All entries are finite by the original Gamma exponential bound stated
below. These Grams are positive definite: their weights
$y^{2(q+s)}\sigma(y)$ are positive away from the single point $y=0$,
and a nonzero polynomial is nonzero on an interval away from that point.
The moment ratio has therefore been specified by complete convergent
integrals; its value is not assumed.

The original full relation Grams and two reference source forms, all on
the same coefficient space $\mathcal P$, are
\[
 \begin{aligned}
 H_{\chi,0}&=\operatorname{Gram}_\sigma
                [\chi S^j]_{j=0}^{r-1},&
 H_{\chi,f}&=\operatorname{Gram}_\sigma
                [f\chi S^j]_{j=0}^{r-1},\\
 H_{\mathrm{pow},s}
   &=\left(\int_{\mathbb R}y^{2(q+s)}
          \overline{(c+iy)^i}(c+iy)^j\sigma(y)\,dy\right)_{i,j=0}^{r-1},
 &s&=0,l .
 \end{aligned}
 \tag{CTR35}
\]
In the first row every polynomial is evaluated at $S=c+iy$.
The coefficient map $P(S)\mapsto P(c+iy)$ is
$J_{h,j}=\binom jh c^{j-h}(\sqrt{-1})^h$ for $h\leq j$,
with zero entries for $h>j$.
Its diagonal entries are $(\sqrt{-1})^h$, so $|\det J|=1$, and
\[
 H_{\mathrm{pow},s}=J^*\mathsf M_{q+s,r}J,\qquad
 \frac{\det H_{\mathrm{pow},l}}{\det H_{\mathrm{pow},0}}
      =\mathcal H_{q,l,r}.
 \tag{CTR36}
\]
This proves the exact translation and phase transport to the real
Hankel matrices, with every coefficient map retained.

Partition the original line as
$I=\{|y|<q\epsilon\}$,
$B=\{q\epsilon\leq|y|\leq64q\}$,
$F=\{|y|>64q\}$. For any one of the four positive full forms in CTR35,
write $H=H^I+H^B+H^F$, and retain the exact positive correction
\[
 t_H=\log\det\left(I_r+
       (H^B)^{-1/2}(H^I+H^F)(H^B)^{-1/2}\right)\geq0,\qquad
 \log\det H=\log\det H^B+t_H.
 \tag{CTR37}
\]
The equality follows by factoring $H$ by the two square roots of $H^B$.
No tail integral has been discarded.

Here are the finite common bounds for these four corrections. Put
\[
 \begin{split}
 \kappa_{I,r}
  &=\frac{2\epsilon r^2}{\sqrt{\cos1}}
       \left(\frac43\,2^{-13}e^{2\pi}\right)^q
    <\frac{2\epsilon r^2}{\sqrt{\cos1}}e^{-2q},\\
 \kappa_F&=\frac{512q}{59\sqrt{\cos1}}e^{-19q},\qquad
 L_r=r\log(1+\kappa_{I,r}+\kappa_F).
 \end{split}
 \tag{CTR38}
\]
For $H_{\chi,0},H_{\chi,f}$ the complete Gamma/Legendre calculation
BRD22--BRD38 proves
$H^I\preceq\kappa_{I,r}H^B$ and
$H^F\preceq\kappa_FH^B$ on exactly the threshold in CTR23.
We verify that the same constants apply to the two explicit reference
forms, rather than transferring a theorem about a different quotient.

The two source bounds proved by the Euler contour rotation and
reflection calculation BRD23--BRD24 are
\[
 C_Le^{-\pi|y|}\leq\sigma(y)\leq C_Ue^{-|y|},\qquad
 C_L=\frac{\pi}{\Gamma(3/4)^2},\quad
 C_U=\frac{\Gamma(1/4)^2}{2\pi\sqrt{\cos1}},\quad
 \frac{C_U}{C_L}=\frac1{\sqrt{\cos1}}.
 \tag{CTR39}
\]
Their full derivation is in the accompanying BRD source. In particular
these are bounds for the same unscaled sigma density as CTR34.
The Rodrigues/orthogonality calculation BRD25--BRD26 gives, for any
polynomial $Q$ of degree at most $r-1\leq q$,
\[
 |Q(z)|^2\leq
 \sum_{j=0}^{r-1}(2j+1)|P_j(2z-3)|^2\int_1^2|Q(v)|^2\,dv.
 \tag{CTR40}
\]
The $P_j$ are the original Legendre polynomials on $[-1,1]$.
For $|z|\leq\epsilon$ their sum is at most $r^2 8^{2q}$;
for $|z|\geq64$ it is at most
$r^2((4+6/64)|z|)^{2r-2}$.
Indeed their elementary real exterior bound is
$|P_j(x)|\leq(2|x|)^j$; substitute $x=2z-3$ and sum
$\sum_{j<r}(2j+1)=r^2$.

For the reference form $H_{\mathrm{pow},s}$, setting $y=qz$ retains
the density factor $q^{2q+2s+1}|z|^{2q+2s}\sigma(qz)$.
On $[1,2]$ its quadratic form is at least
\[
 q^{2q+2s+1}C_Le^{-2\pi q}\int_1^2|Q(z)|^2\,dz.
 \tag{CTR41}
\]
On the inner region the numerator is bounded by
$q^{2q+2s+1}\epsilon^{2q+2s}C_Ur^28^{2q}$
times $2\epsilon\int_1^2|Q|^2$. Dividing by CTR41 proves
\[
 H_{\mathrm{pow},s}^I\preceq
 \frac{2\epsilon r^2}{\sqrt{\cos1}}\,
       \epsilon^{2s}\bigl(2^{-14}e^{2\pi}\bigr)^q
       H_{\mathrm{pow},s}^B
 \preceq\kappa_{I,r}H_{\mathrm{pow},s}^B.
 \tag{CTR42}
\]
Here $\epsilon=2^{-10}$ makes $\epsilon^2 8^2=2^{-14}$,
$\epsilon^{2s}\leq1$, and
$2^{-14}<(4/3)2^{-13}$. This is the reference's actual
inner integral estimate at zero root radii.

For its far region, put $D_s=2q+2s+2r-2\leq5q$.
The source and polynomial bounds in CTR39--CTR40, divided by CTR41,
give
\[
 H_{\mathrm{pow},s}^F\preceq
 \frac{2r^2}{\sqrt{\cos1}}\,
 \frac{e^{2\pi q}(131/32)^{2r-2}64^{D_s}e^{-64q}}
      {q-D_s/64}\ H_{\mathrm{pow},s}^B.
 \tag{CTR43}
\]
To verify the integration constant, for $z\geq64$ use
$\log(z/64)\leq(z-64)/64$, which gives
$\int_{64}^\infty z^{D_s}e^{-qz}dz
\leq64^{D_s}e^{-64q}/(q-D_s/64)$.
The factor two counts both real tails. Since
$q-D_s/64\geq59q/64$ and
$64-2\pi-5\log64-2\log(131/32)>20$, the coefficient in CTR43 is
at most
$128r^2e^{-20q}/(59q\sqrt{\cos1})\leq\kappa_F$,
using $r\leq q+1\leq2q$.
Thus all four actual corrections in CTR37 obey
\[
 0\leq t_H\leq L_r.
 \tag{CTR44}
\]
This follows by congruence, since every eigenvalue of its positive
tail matrix is at most $\kappa_{I,r}+\kappa_F$.

On the common bulk, preserve $f$ while applying CTR28. Define
\[
 E_{\chi,l}
   =l\{|a_k|(\epsilon^{-2}-64^{-2})+2\delta_\chi\},\qquad
 \Delta_f=\frac{S_2}{q^2\epsilon^2},\qquad
 S_2=\sum_{j=0}^{l-1}(2j+\tfrac12)^2
     =\frac{l(16l^2-12l-1)}{12}.
 \tag{CTR45}
\]
The signed bulk source difference
$d_\chi=\log\mathcal R_{w_\chi}(f)
       -\log\mathcal R_{w_{\mathrm{pow}}}(f)$ satisfies
$|d_\chi|\leq E_{\chi,l}$.
For the remaining multiplication factor, at every bulk point
\[
 1\leq\frac{|f(c+iy)|^2}{|y|^{2l}}
       =\prod_{j<l}(1+t_j^2/y^2)\leq e^{\Delta_f}.
 \tag{CTR46}
\]
The upper bound follows by $\log(1+x)\leq x$ and $|y|\geq q\epsilon$.
Integrating every squared polynomial gives the positive-form comparison
between the bulk numerator with $f$ and the bulk reference numerator
with $|y|^{2l}$. Congruence and multiplication of its $r$ eigenvalues
therefore give a one-sided bulk correction $d_f$ satisfying
$0\leq d_f\leq r\Delta_f$. Its denominator is the identical
$H_{\mathrm{pow},0}^B$.

Writing the full original relation determinant as
$\mathfrak R_{q+r-1}=\det H_{\chi,f}/\det H_{\chi,0}$,
CTR36--CTR37 now prove the complete exact identity
\[
 \begin{split}
 \log\mathfrak R_{q+r-1}-\log\mathcal H_{q,l,r}
   ={}&d_\chi+d_f
       +t_{H_{\chi,f}}-t_{H_{\chi,0}}\\
      &-t_{H_{\mathrm{pow},l}}+t_{H_{\mathrm{pow},0}} .
 \end{split}
 \tag{CTR47}
\]
All four original or reference nonnegative tail corrections remain
separately in this formula. In particular its proved finite enclosure is
\[
 \boxed{
 -E_{\chi,l}-2L_r
 \leq\log\mathfrak R_{q+r-1}
       -\log\frac{\det\mathsf M_{q+l,r}}{\det\mathsf M_{q,r}}
 \leq E_{\chi,l}+r\Delta_f+2L_r,\qquad r=q,q+1.}
 \tag{CTR48}
\]
The one-sided $r\Delta_f$ appears only in the upper endpoint.
Equations CTR31--CTR33 bound $E_{\chi,l}=O(k/e)$.
Also $S_2\leq4l^3\leq k^3/16$ and $r\leq2q$ imply
$r\Delta_f\leq k^3/(8q\epsilon^2)\leq k/(8e\epsilon^2)$.
Finally $L_r\leq r(\kappa_{I,r}+\kappa_F)$ decays exponentially
in $q$ times an explicit polynomial by CTR38.
It follows from these displayed constants that
\[
 \log\mathfrak R_{q+r-1}
   =\log\frac{\det\mathsf M_{q+l,r}}{\det\mathsf M_{q,r}}
       +O(k)=
     \log\frac{\det\mathsf M_{q+l,r}}{\det\mathsf M_{q,r}}+o(q)
 \tag{CTR49}
\]
for each fixed original packet and $m$, and the error is $O(1)$
when the fixed multiplicity satisfies $m\geq2$.
The reference now depends only on the explicit integers $q,l,r$;
$\delta,\gamma$ remain in the proved error constants and threshold.

\subsection{The full reference entries are exact sigma moments}
To give a complete specification without unevaluated moment inputs,
write $\mu_j=\int_{\mathbb R}y^j\sigma(y)\,dy$.
The characteristic function JSA4 of the same original sigma density is
\[
 \int_{\mathbb R}e^{ity}\sigma(y)\,dy
    =c_\sigma(\cosh t)^{-1/2},\qquad c_\sigma=\sqrt{2\pi}.
 \tag{CTR50}
\]
Its direct derivation can be retained here. In the product of the Euler
integrals for $\Gamma(a+iy/2)$ and $\Gamma(a-iy/2)$, use
$u+v=w$ and $\log(u/v)=2x$. The Jacobian and the positive $w$ integral
give
$|\Gamma(a+iy/2)|^2
=2\Gamma(2a)\int_{\mathbb R}e^{iyx}(2\cosh x)^{-2a}dx$.
Fourier inversion gives
$2\Gamma(2a)(2\cosh t)^{-2a}$ for its density characteristic function.
The function $(2\cosh x)^{-2a}$ is smooth with integrable exponentially
decaying derivatives; Fourier inversion also follows by Gaussian
regularization and dominated convergence. At $a=1/4$ the constant is
$2\Gamma(1/2)2^{-1/2}=\sqrt{2\pi}$, proving CTR50.

The exponential upper bound in CTR39 permits differentiation of every
finite order under the integral in CTR50: the derivative is dominated
by $|y|^jC_Ue^{-|y|}$. Thus
\[
 \mu_{2h+1}=0,\qquad
 \mu_{2h}
  =c_\sigma(-1)^h
     \left.\frac{d^{2h}}{dt^{2h}}(\cosh t)^{-1/2}\right|_{t=0},
 \qquad
 (\mathsf M_{a,r})_{i,j}=\mu_{2a+i+j}.
 \tag{CTR51}
\]
The odd moments vanish either by the even characteristic function or
by the original even density. These formulas retain the unscaled
variable $y$ and the full Gamma mass.

For an explicit finite rational recurrence let
$c_h=((d/dt)^{2h}(\cosh t)^{-1/2})|_{t=0}$ and $c_0=1$.
The derivative equation
$h'(t)\cosh t+\tfrac12h(t)\sinh t=0$ for $h(t)=(\cosh t)^{-1/2}$
gives, by comparing its coefficient of $t^{2n-1}$,
\[
 c_n=
 -\sum_{j=1}^{n-1}\binom{2n-1}{2j-1}c_j
 -\frac12\sum_{j=0}^{n-1}\binom{2n-1}{2j}c_j
 \quad(n\geq1).
 \tag{CTR52}
\]
In this comparison the coefficient of $c_n$ is one; the remaining
terms come respectively from $h'\cosh$ and $h\sinh/2$.
It proves that every $c_n$ is rational and determines it in finitely
many operations. For example $c_1=-1/2$ and $c_2=7/4$.
Hence $\mathsf M_{a,r}=c_\sigma\mathsf Q_{a,r}$, where the rational
matrix has entry zero at odd $i+j$, and entry
$(-1)^{a+(i+j)/2}c_{a+(i+j)/2}$ at even $i+j$.
In particular
\[
 \det\mathsf M_{a,r}=c_\sigma^r\det\mathsf Q_{a,r}>0,\qquad
 \mathcal H_{q,l,r}
   =\frac{c_\sigma^r\det\mathsf Q_{q+l,r}}
          {c_\sigma^r\det\mathsf Q_{q,r}} .
 \tag{CTR53}
\]
The mass factors are displayed in both determinants before cancellation.
Thus the new reference is an exact positive rational Hankel-determinant
ratio, supplied by a finite recurrence, rather than an unspecified
comparison measure.

The existing low-endpoint constant $C_{\rm rel}$ is retained unchanged
in the original identity
$\mathfrak T_k=C_{\rm rel}+\log\mathfrak R_{2q-1}
                              +\log\mathfrak R_{2q}$.
Substitution of CTR47 gives
\[
 \mathfrak T_k=C_{\rm rel}
      +\sum_{r=q}^{q+1}
          \log\frac{\det\mathsf M_{q+l,r}}{\det\mathsf M_{q,r}}
      +\varepsilon_q+\varepsilon_{q+1},
 \tag{CTR54}
\]
where each $\varepsilon_r$ has the exact four-tail decomposition
CTR47 and the one-sided bounds CTR48. Their sum is $O(k)=o(q)$,
and $O(1)$ for fixed $m\geq2$. The incoming Gamma return is
$R_k^0-R_k^\sigma=-\mathfrak T_k$, so its entire expression has the
opposite sign. The next calculation is the asymptotic evaluation of
these specified universal Hankel determinants; no asymptotic value
has been inserted in CTR49 or CTR54.
